% !TEX root = ../main.tex

\chapter*{出版说明}
\addcontentsline{toc}{chapter}{出版说明}

本书为《杂阿含经·仿罗什风新译》之\textbf{法义读本}（V3）。
编辑假想是：\textbf{若鸠摩罗什为通读而彻底重译阿含}，
必不按大正经目短题逐条照录，而会重拟篇章、删并重复、一气连贯——
如其译《金刚经》之删冗圆通。

\vspace{0.8em}
\noindent\textbf{可读性第一：}
\begin{itemize}
  \item 篇题均为\textbf{通读新拟}（如「五阴无常」「十二支流转」），\textbf{不是}「如来第一」「离贪法第一」一类原经短题。
  \item 每篇熔多经罗什风译文为一篇连续正文，去掉反复的「如是我闻……欢喜奉行」。
  \item 只保留【正文】与【今译意】；附注说明所熔经号，供回指 V1／V2。
\end{itemize}

\vspace{0.8em}
\noindent\textbf{与研究本：}V1／V2 仍是逐经全收、可核验的研究译注本；本书不取代它们，也不伪托「罗什历史译本」或「原典次序」。

\vspace{0.8em}
\noindent 原则论证见 \texttt{docs/V3\_DHARMA\_READER.md}。

\clearpage
